\section{Conclusiones}

El proceso de análisis del Modelo Conceptual para el sistema ``Gemelo Digital EAM'' ha permitido consolidar una visión técnica de alta fidelidad que trasciende la simple digitalización, estableciendo un marco de gobernanza operativa basado en datos maestros. La implementación del \textbf{Sidecar Pattern} y la arquitectura de \textbf{Requisitos Transversales (TR)} asegura una base desacoplada y escalable, permitiendo que la visualización 3D evolucione sin comprometer la integridad del historial de mantenimiento.

La alineación estricta con las normas \textbf{ISO 55001} e \textbf{ISO 14224} dota al sistema de un rigor industrial necesario para la toma de decisiones basada en la confiabilidad real de los activos. Al actuar como un instrumento de auditoría y mejora continua, la solución se posiciona como una herramienta estratégica para la gestión de activos críticos en el sector industrial.

El diseño de la lógica de aislamiento bajo procedimientos \textbf{LOTO} y criterios de integridad \textbf{fail-safe} demuestra un compromiso con la seguridad funcional y la protección del personal operativo. Esta integración eleva el software de una plataforma administrativa a un sistema de seguridad técnica esencial para la operación en planta.

La incorporación de modelos de asistencia basados en \textbf{RAG} y la sincronización con el estándar \textbf{HPHMI} resuelve problemáticas históricas de carga administrativa y fatiga visual en campo. Estas innovaciones permiten que el personal especializado optimice su tiempo en actividades de alto valor, mientras que el enfoque \textbf{Offline-First} garantiza la resiliencia operativa y la trazabilidad de los datos en cualquier condición de conectividad de la planta.

Los 360 requisitos refinados y modelados en este documento constituyen la hoja de ruta técnica para la fase de implementación, asegurando el cumplimiento de los objetivos estratégicos de la organización y los estándares de calidad del SENA.
